% =====================================================================
%  第 1 章 绪论 —— Lead 亲译的样板章（ch01）
%  页段：PDF 11–30（原书 p.1–20）
% =====================================================================
\chapter{绪论}

% [原文 p.1]
现在写一本辛几何的教科书还为时尚早。

辛几何发端于 19 世纪数学家 Hamilton：它作为把古典力学方程组的变分法结构表达出来的数学型式而登场。

奇妙的是，本来应当是古典力学表达手段的辛几何（哈密顿形式），在量子力学的创立中也扮演了不可或缺的角色。

在 20 世纪几何学的发展中，辛几何始终是一个异类。在几何学“从局部到整体”的潮流中，辛几何完全落了伍。因此，它虽作为未来数学的萌芽被一部分人热烈拥护，并在几何学以外的领域催生了许多重要应用，却终究没有成为几何学本身的主流。

Gromov 1985 年的论文终于揭示出整体辛几何具有丰富的内容，一举改变了那幅图景。当理论物理学中整体几何学的意义终于被发现时，整体辛几何的重要性就愈发显而易见了。

辛几何的时代，现在才刚刚开始。

\section{埃尔朗根纲领与 G 结构}

% [原文 p.2]
开头那段引文，是笔者为岩波讲座《现代数学的展开》所写的内容预告。我想借发挥它来讲述辛几何的历史——一面注视 20 世纪后半叶几何学“从局部走向整体”的潮流，一面回顾 20 世纪几何学走到这一步的历程。

\takeaway{本节把“几何”统一为群作用下的不变量理论：埃尔朗根纲领用对 $(X,G)$ 刻画几何学，Cartan 的 G 结构则把它局部化到切空间 $T_pM$。可积性把近复结构（$GL(n;\C)$ 结构）与复结构区分开；辛结构正是 $G=Sp(n)\subset GL(2n;\R)$ 的 G 结构，其可积性条件为 $d\omega=0$。}

就从著名的 Klein 埃尔朗根纲领（Erlangen program）讲起吧。先简化地表述一下埃尔朗根纲领：给定群与它作用的空间组成的对 $(X,G)$，就给出一种几何学。所谓具有与 $(X,G)$ 相对应的几何学的空间 $M$，是指：对任意点 $p\in M$，其邻域 $U_p\subset M$ 都可以视为 $X$ 的开集；并且，$X$ 上的诸量或诸概念中，只有那些在 $G$ 作用下不变者，才被视为由 $(X,G)$ 给出的几何学的概念。

\begin{remark}{1.1}
  若精确地定义，则如下所述。

  对 $M$ 的每一点 $p$，若给定其邻域 $U_p$ 与 $X$ 的开集之间的微分同胚 $\varphi_p$，使下述条件成立，则称 $M$ 上确定了一个 $(X,G)$ 结构：坐标变换 $\varphi_p\circ\varphi_q^{-1}\colon\varphi_q(U_p\cap U_q)\to\varphi_p(U_p\cap U_q)$ 可用 $G$ 的元素 $g_{p,q}$ 表示为 $\varphi_p\circ\varphi_q^{-1}(x)=g_{p,q}x$。
\end{remark}

% [Tutor Note: 这里的 $g_{p,q}$ 依赖于点对 $(p,q)$，不是常值矩阵；G 结构的“微分”由 $G$ 给出，正是指这个逐点取自 $G$ 的过渡函数。$x\in U_p\cap U_q$ 变化时 $g_{p,q}$ 随之变化，这正是“G 结构未必是平凡化”的要点。]

% [原文 p.3]
典型的例子是：$X$ 为 $n$ 维欧氏空间、$G$ 为其合同变换（等距变换）全体之群的情形（欧氏几何）；$X=\R^n$、$G$ 为其相似变换全体之群的情形（相似几何）；$X$ 为曲率 $-1$ 的常曲率空间、$G$ 为等距变换之群的情形（双曲几何）；$X$ 为射影空间 $\CP^n$、$G$ 为 $GL(n+1;\C)$ 的情形（射影几何）；等等。

% [Tutor Note: 对传递作用的情形，空间在群意义下就是齐性空间 $G/H$（$H$ 为一点稳定子群）。埃尔朗根纲领与 Cartan 的 G 结构的区别在于：前者要求“有限邻域”与 $X$ 的开集等同，后者只要求“切空间”等同。]


这里并不打算逐一讲解，故从略。笔者想说的是这样一种思想：给定群，就得到一种几何学。在埃尔朗根纲领中，“给定群”就是指给定 $(X,G)$ 这样的对。群 $G$ 作用在空间 $X$ 上。在重要的例子里，$G$ 对 $X$ 的作用是传递的（transitive）。（也就是说，当 $p,q\in X$ 时，存在 $g\in G$ 使 $q=gp$。）此时，所考虑的空间大致是齐性空间 $G/H$。（这里 $H$ 是 $G$ 的子群，$G/H$ 指 $H$ 从右侧作用于 $G$ 所得的商空间。）

把这一框架局部化的人是 E. Cartan。在 Cartan 的几何学中，群 $G$ 作用在向量空间 $\R^n$ 上，即 $G\subset GL(n;\R)$。而 $n$ 维流形 $M$ 的 \textbf{G 结构}（$G$ structure）是指 $M$ 的坐标系 $(U_i,\varphi_i)$，其坐标变换的“微分”由 $G$ 的作用给出。

也就是说，在埃尔朗根纲领的几何学中，空间 $M$ 在各点的邻域处被与 $(X,G)$ 的一部分等同看待；与此相对，在 Cartan 的几何学中，$M$ 各点的无穷小邻域即切空间 $T_pM$ 被与 $(\R^n,G)$ 等同看待。

\begin{remark}{1.2}
  精确地说，是指：存在 $M$ 的开覆盖 $M=\bigcup U_i$，以及微分同胚 $\varphi_i\colon U_i\to\varphi_i(U_i)$（$\varphi_i(U_i)$ 为 $\R^n$ 的开集），使得 $\varphi_j\circ\varphi_i^{-1}\colon\varphi_i(U_i\cap U_j)\to\varphi_j(U_i\cap U_j)$ 在 $x\in\varphi_i(U_i\cap U_j)$ 处的微分
  \[
    d(\varphi_j\circ\varphi_i^{-1})\colon T_x\R^n \longrightarrow T_p\R^n
  \]
  是 $G\subset GL(n;\R)$ 的元素（其中 $p=\varphi_i^{-1}(x)$）。
\end{remark}

G 结构的两个典型例子是：黎曼几何（Riemannian geometry），即 $G=O(n)\subset GL(n;\R)$（正交群）的情形；以及复几何，即 $G=GL(n;\C)\subset GL(2n;\R)$ 的情形。

把埃尔朗根纲领与 Cartan 的几何学联系起来的，是 G 结构的\textbf{可积性}（integrability）这一概念。G 结构可积，定义如下。当 $G\subset GL(n;\R)$ 时，$\R^n$ 上确定有标准的 G 结构。$M$ 上的 G 结构可积，是指：对 $M$ 的每一点，其某个邻域连同 G 结构一起，同构于带有标准 G 结构的 $\R^n$ 的开集。

在 G 结构的定义中，各点的无穷小邻域被与带有标准 G 结构的 $\R^n$ 等同看待。而可积性意味着：一个充分小、但并非 0 的有限大小的邻域，可以与带有标准 G 结构的 $\R^n$ 等同看待。

% [Tutor Note: “可积”在这里是几何意义而非微分方程意义：$G$ 结构可积 $\Leftrightarrow$ 存在局部坐标系使结构张量成为常值标准形。对辛结构（$G=Sp(n)$）而言，可积性由 $d\omega=0$ 保证（Darboux 定理，§2.2）；对近复结构则是 Newlander–Nirenberg 定理。]

复结构就是可积的 $GL(n;\C)$ 结构。在不假定可积性时，$GL(n;\C)$ 结构称为\textbf{近复结构}（almost complex structure）。
% [Tutor Note: 日语“概複素構造”字面为“almost complex structure”，中文通行译名为“近复结构”。同理本书的“概正則曲線”译作“伪全纯曲线”。全书统一。]

% [原文 p.4]
称为近复结构。（这就是上面在“即”字上加括号的理由。）

在黎曼几何中，并不假定可积性。在黎曼几何的情形，即 $G$ 为正交群 $O(n)\subset GL(n;\R)$ 的情形，可积的 G 结构意味着平坦的黎曼度量；或者说，意味着 $X$ 为 $n$ 维欧氏空间、$G$ 为其合同变换（等距变换）之群时的 $(X,G)$ 结构。这作为研究对象未免太窄了。

Cartan 的几何学把不可积的 G 结构也纳入进来，从而把几何学的对象扩展到齐性空间之外的一般流形。齐性空间的研究中，李群或李代数的代数研究曾是最有力的手段。超越齐性空间，就要求超越代数方法、让超越的方法成为研究的中心手段。然而，正如 Cartan 本人就是李群论的主要奠基人之一这一点所表明的，超越的方法要成为几何学的主要方法，还需历经许多年月。

20 世纪前半叶微分几何学的发展，是以局部理论为中心展开的。用 G 结构的语言来说，这就是具有 G 结构的空间的局部分类问题。只要考虑可积的 G 结构，局部分类问题就不存在。因为“局部上与标准者相同”正是可积性的定义。但若考察黎曼几何，局部黎曼流形的分类问题，就是自研究曲面论的 Gauss 以来的中心问题。

给定 G 结构，多数情况下意味着在空间上给定一个决定结构的张量。例如在黎曼几何中，它是正定对称张量 $g\colon TM\otimes TM\to\R$；而在（近）复结构中，它是相当于乘以 $\sqrt{-1}$ 的映射（1-1 张量）$J\colon TM\to TM$。

要回答 G 结构的局部分类问题，必须找到 G 结构的不变量。它由对给定 G 结构的张量求导所得的量适当组合而得到。作怎样的组合，所得之物才是 G 结构的不变量，即在保持 G 结构的微分同胚下保持不变——这就是基本问题。这样的量称为微分不变量（differential invariant）。

% [原文 p.5]
求微分不变量的问题，自 19 世纪末起就作为张量分析的中心问题被研究。在黎曼几何的情形，曲率是最重要的微分不变量。

G 结构的局部分类问题，只要找到足够多的微分不变量，并反过来证明：这些微分不变量一致的两种 G 结构局部同构，就算解决了。这称为等价问题（equivalence problem）。

在黎曼几何的情形，与之相当的命题是“曲率一致的两个黎曼流形局部等距”。（若不给出“曲率一致”的准确含义，这个命题就不构成数学定理。此处从略。参见 [204] Vol.1, Chapter IV 7 等。）

等价问题的主要部分到 1950 年代已获解决，微分几何学与拓扑几何学的基础也告完成。于是微分几何学的中心转向整体问题。

\section{从局部到整体}

辛几何的中兴之祖 Gromov，1985 年在东京大学就辛几何作了演讲。但当时日本几乎没有研究 Gromov 所展开的那种类型辛几何的研究者，因而并没有出现现场热烈讨论的场面。演讲结束后，Gromov 转向笔者，不满地问：为什么没有人提问？笔者就 Gromov 所提出的辛几何，询问其目标与动机。对此，Gromov 简单地答了一句：

辛几何比黎曼几何更丰富。

这或许也是意识到当时正在研究黎曼几何的笔者而说的话\footnote{黎曼几何也与辛几何同样丰富。}；如今想来，这也许正表达了他对自己所提出的、作为“通常的整体几何学”之又一新成员的辛几何的可能性所抱的确信。

\takeaway{Gromov 把辛几何定位于“通常的整体几何学”：几何学的中心从局部不变量转向整体问题后，辛几何要问的正是存在性与分类。这一转向以 Gromov 1985 年的伪全纯曲线方法为标志，也是全书第 2 章以后的主线。}

% [原文 p.6]
刚才我写了“通常的整体几何学”这个词（这是笔者随意使用的）。我想说明它的含义。

20 世纪后半叶，整体研究跃居几何学的前台。在 1950 年代承担这一转向的动向，是微分拓扑学的创立，以及复几何学从代数几何学与多复变函数论中独立出来。

先前讨论几何结构时，我们一开始就把微分流形的结构当作给定之物。也就是说，把各种几何结构（G 结构）看作微分结构的下层结构。若把出发点不放在微分流形上而放在拓扑流形上，则可以把微分结构本身看作一种几何结构。这时“考虑群 $G\subset GL(n;\R)$”这一视角需要稍作修正，但也仍然可行。（考虑欧氏空间局部微分同胚之群 $\Diff$ 与局部同胚之群 $\Top$，用 $\Diff\subset\Top$ 这一对替换 $G\subset GL(n;\R)$。此时，切丛（tangent bundle）的概念可推广到拓扑流形，这一点在 60 年代由 Milnor 等人以微丛（micro bundle）理论建立起来。在拓扑流形的微丛中，$\Top$ 扮演向量丛中 $GL(n;\R)$ 的角色。）

把微分结构看作“几何结构”时，它当然是“可积”的。也就是说，对微分流形 $M$，点 $p\in M$ 的邻域与欧氏空间的开集微分同胚。因此，在微分拓扑学中，分类的局部问题是平凡的。然而，分类的整体问题，即“两个微分流形 $M,M'$ 同胚时是否微分同胚”，或者“对给定的拓扑流形，是否存在与之同胚的微分流形”，却是深刻的问题。

笔者不知道在 Milnor 给出否定解答之前，这些问题是否已被意识到是问题。Milnor 对第一个问题在 7 维球面的情形给出了反例，Kervaire 则对第二个问题给出了反例。而且这些反例并非病态的例子，而是普通地……

% [原文 p.7]
……只要按通常方式考虑，就是自然出现的正统对象。

就这样，上述两个问题被公认为微分拓扑学的基本问题，微分拓扑学作为数学的一个领域正式开始了。

从同一处出发，整体黎曼几何也逐步发展起来。整体黎曼几何的基本问题是：曲率对整体的性质有何影响？

微分拓扑学的发展，逐步弄清了存在哪些微分流形、要区分它们需要考察什么，从而为大范围黎曼几何准备了必要的骨架。另一方面，取代作为局部黎曼几何主要手段的张量分析，整体黎曼几何的方法也开始出现。其一是关于曲线长度的变分法即 Morse 理论；另一是调和积分论（harmonic analysis，由 Hodge 等人建立）与 Bochner 技巧（Bochner trick），即曲率与流形上偏微分方程之间的关联。

% [Tutor Note: Morse 理论用变分法（常微分方程层面）提取拓扑信息；调和积分与 Bochner 技巧则把曲率条件转化为偏微分方程（椭圆型）的可解性。这三条线索（变分、椭圆方程、非线性方程）正是整体几何学的三大工具箱。]


另一方面，复几何的独立，很大程度上归功于超越的方法、特别是使用偏微分方程的方法在复流形上的应用得以确立。这是由 Hodge、De Rham、小平等人建立的调和积分论，以及 Gauss–Bonnet 定理（的高维化）、Riemann–Roch 定理、Atiyah–Singer 指标定理这一系列研究所成就的；小平在这些工作的基础上，把复几何确立为一个领域。复流形的同型类的个数（不同于微分流形的微分同胚类的个数）是不可数的。因此，“分类”离不开复流形同型类所成模空间（moduli space）的构造与研究。小平与 Spencer 一起开创了复流形的形变理论（deformation theory），即模空间的局部理论。另一方面，小平同样超越了古典已知的复 1 维情形（即黎曼面的情形），深入研究复 2 维复流形的分类即复解析曲面的分类理论，补足了意大利学派代数曲面论的理论不足。

微分流形（5 维以上）的分类、复流形模空间的理论、复曲面论在 50、60 年代的发展，使整体几何学成为几何学（或者说数学）的中心。

% [原文 p.8]
到了 70 年代，整体几何学又添了一种重要手法，即非线性偏微分方程在几何学中的应用。它与前面作为整体黎曼几何两种方法所述的、关于曲线长度的变分法（这是非线性常微分方程的问题）以及线性偏微分方程相并列，使流形上的整体微分几何在 70 年代以后取得了飞跃进展，许多未解决问题获得解决。其基本问题意识，与微分拓扑学诞生处所述者类似。也就是说，考察满足给定条件的黎曼流形或复流形是否存在，或者对它们进行分类。这样的数学，这里权且称为“通常的整体几何学”。

在 Gromov 的论文 [159] 出现之前，辛几何走的似乎是与此不同的道路。

\begin{remark}{1.3}
  作为撰写辛几何史的人，笔者完全不合适\footnote{笔者所写的第一篇关于辛几何的论文是 1993 年前后的，在 Gromov 以后的辛几何史中也属较晚的。}。这里所写的，深受身为“通常的整体几何学”研究者的笔者在 1980 年前后对辛几何所抱印象的影响。例如，从一直把辛几何作为一大潮流加以研究的俄罗斯数学家看来，或许会给出完全不同的看法。
\end{remark}

把辛结构用 G 结构的语言写出来，就可以说是 $G=Sp(n)\subset GL(2n;\R)$ 的情形。这里
\[
  Sp(n)=\left\{A\in GL(2n;\R)\ \middle|\
  A^{-1}\begin{pmatrix}0&-I_n\\ I_n&0\end{pmatrix}A
  =\begin{pmatrix}0&-I_n\\ I_n&0\end{pmatrix}\right\},
\]
（$I_n$ 为单位 $n\times n$ 矩阵）。换言之，$M$ 的“辛结构”就是非退化反对称 2 次形式 $\omega\colon TM\times TM\to\R$。（黎曼度量则是正定对称 2 次形式。）反对称 2 次形式 $\omega\colon TM\times TM\to\R$ 也就是微分 2-形式，所以“辛结构”就是非退化微分 2-形式。
% [Tutor Note: 这里把辛结构写成 $G=Sp(n)$ 结构，是全书反复使用的框架：辛形式 $\omega$ 在每点给出一个标准反对称形式，$Sp(n)$ 正是保持该形式的线性群。注意与后面 §2.3 的“近复结构 $J$”配合时，$Sp(n)\cap O(2n)=U(n)$，这是凯勒流形的群论刻画。]

% [原文 p.9]
上面加括号，是因为辛结构还需要另一个条件，即可积性条件。在辛结构的情形，可积性条件等价于 $d\omega=0$ 这一条件。（Darboux 定理。§2.2 详述。）因此，仅仅非退化的微分 2-形式，若模仿复结构的情形，或许该称为近辛结构。（至于近辛结构的局部理论能否像黎曼几何那样具有丰富的内容，笔者不得而知。）

在辛几何的情形，也可以考虑上一节所述“通常的整体几何学”的问题。然而，基于这种问题意识的发展，长期以来似乎并不是辛几何的中心。这些问题的解答开始被发现、作为“通常的整体几何学”的辛几何开始发展，是在 Gromov 的论文 [159] 之后。关于这一发展，将在第 2 章以后详述，此处不涉及。此处倒想谈谈此前的辛几何。

\begin{remark}{1.4}
  说“此前的辛几何”这种措辞，听起来总有些陈旧。但事实并非如此。辛几何长期以来怀着对 21 世纪的梦想——不，莫如说是在梦想几乎要把自己压垮的情形下——发展起来的。它把量子化树为目标，却并不像“通常的整体几何学”那样，拥有许多可以拿出来说“这就是成果”的、漂亮的、数学上明确的定理。正因如此，笔者如今才对这一潮流格外在意。
\end{remark}

至此所述的 20 世纪几何学史，是正统的看法，即正史。当 21 世纪几何学诞生之时，正史中被轻视的其他研究、在华丽发展之外流淌的其他潮流，或许会浮现到前台。若辛几何成为孕育 21 世纪几何学的一大潮流……

% [原文 p.10]
……那么就必须找出正史中已记载的 Gromov\footnote{本节与后面各节的说明，写法上都有可能招致“大范围辛几何是否仅由 Gromov 的论文 [159] 开始”的误解。[159] 的重要性是显然的，但在此之前为孕育大范围辛几何做出贡献的许多人当然存在。他们的名字大多将在后面的章节中引用，此处不举，以 [159] 代表大范围辛几何创生期许多人的努力成果。} 之后的研究与此前研究之间的接点。

\section{辛几何历史一瞥}

我的想法是：全部数学的辛几何化。

这是长期支撑辛几何发展的 Arnold 写给笔者的电子邮件（1996 年）中的一段。Arnold 无疑也是梦想着以辛几何开创 21 世纪几何学的人之一。

\takeaway{本节从 Hamilton 与 Lagrange 两种力学的对等性出发，说明辛形式是 Hamilton 力学的几何内核；又借 Arnold 与作者之口指出，辛几何的历史遗产在于“量子化”这一未竟目标，正史之外的潮流正在为 21 世纪的几何学积蓄力量。}

辛几何的历史很古老。它的创始者，是建立分析力学哈密顿形式（Hamilton formalism）的 Hamilton。让我们回想哈密顿形式与拉格朗日形式（Lagrange formalism）的区别。（详见例如 [13]、[127]。）对于由在 3 维欧氏空间中运动的 $n$ 个粒子组成的系统，粒子的运动由映射 $q\colon\R\to\R^{3n}$ 表示。在拉格朗日形式中，考虑 $q$ 的（泛）函数 $L(q)$，求其变分问题。与此相对，在哈密顿形式中，把表示位置的映射 $q$ 与表示动量的映射 $p\colon\R\to\R^{3n}$ 都看作自变量，再应用变分法。在拉格朗日形式中，拉格朗日泛函 $L(q)$ 当然也含有 $q$ 的微分项。但在拉格朗日形式中，$q$ 的微分被认为由 $q$ 自动决定，不看作自变量。而在哈密顿形式中，$p=m\,dq/dt$ 这一式子，是作为 $(q,p)$ 给出哈密顿泛函极值这一事实的推论之一而出现的。

% [Tutor Note: 把 $p=m\,dq/dt$ 视为“推论”而非“定义”，是 Hamilton 形式与 Lagrange 形式的本质差异：前者把 $q,p$ 同列为自变量，后者只以 $q$ 为自变量。用几何的语言，这对应“切丛 $TM$ 上的二阶方程”与“余切丛 $T^*M$ 上的辛结构”两种描述。]


这一差别是重要的。也就是说，在哈密顿形式中，$q$ 与 $p$ 的角色是对等的。这一点意味着，哈密顿形式比拉格朗日形式容许更多的变换……

% [原文 p.11]
……也就是说，在拉格朗日形式中，只允许 $\R^n$ 的坐标变换、即把 $q$ 变为 $q$ 的变换；而在哈密顿形式中，则允许把 $q,p$ 任一作为自变量的变换。不过并非允许所有变换，只有保持微分形式
\begin{equation}
  \omega=dq_1\wedge dp_1+\cdots+dq_{3n}\wedge dp_{3n} \tag{1.1}
\end{equation}
的变换，才能使哈密顿形式不变。这个 $\omega$ 就是辛形式，辛几何的起源正在于此。哈密顿形式允许更大的变换群，其好处首先在于：利用更大的变换自由度，把方程化为简单形式的可能性变大了；但意义不止于此。

让我们考虑并非欧氏空间的一般空间、例如流形 $M$ 上的运动。此时可以把 $q$ 看作 $M$ 上的函数。$p$ 又是什么呢？仔细考虑式 (1.1) 的坐标变换性便可明白：把 $(q,p)$ 这一对看作到余切丛 $T^*M$ 的映射是恰当的。换言之，$T^*M$ 上自然地带有辛结构（参见 §2.1 的引理 2.2）。

然而，一旦这样看，前面见到的那个重要对称性、即 $p$ 与 $q$ 地位对等，就被破坏了。也就是说，$q$ 一般是弯曲空间 $M$ 上的坐标函数，而 $p$ 则是作为余切丛 $T^*M$ 的纤维的向量空间上的坐标函数。

要恢复 $p$ 与 $q$ 的对称性，最简便的办法是把视野扩大到不限于余切丛 $T^*M$ 的一般的辛流形。也就是说，把一般流形 $X$ 与其上满足 $d\omega=0$ 的非退化微分 2-形式 $\omega$ 组成的对 $(X,\omega)$，看作展开“流形上的分析力学”的舞台。这暂且就是辛几何的立场。

这样想时，会打开一片诱人的视野，即：把辛几何看作通常几何学的推广。下面来说明。

考虑一般的辛流形 $X$，若用 $X$ 上的函数即哈密顿泛函作变分法就是哈密顿形式，那么与之对应的拉格朗日形式是什么呢？辛流形 $X$ 是……

% [原文 p.12]
……当 $X$ 是余切丛 $T^*M$ 时，考虑关于到 $M$ 的映射 $q\colon\R\to M$ 的泛函，就能找出等价的拉格朗日形式。然而，当 $X$ 不是余切丛 $T^*M$、例如 $X$ 紧致时又当如何？什么是 $M$？或者，这个哈密顿形式描述的是哪个空间上的运动？

对这个提问，不妨这样回答。

当 $X$ 不是余切丛 $T^*M$ 时，对应的 $M$ 在通常意义下（例如作为流形）不存在。也就是说，一般 $X$ 上的哈密顿形式，可视为对应着描述属于某种未知几何学的“空间”上运动的拉格朗日形式。或者也可以说，$M$ 上的辛几何，与属于新空间概念的某个空间上的几何等价。

照这样，这个回答过于轻易，似乎不会由此立即引出很多东西。然而，梦想着以辛几何开创 21 世纪几何学这一立场，最朴素的表达正在于此。

故事并未到此结束。我们只是看到了辛流形概念的出现，恐怕是在 19 世纪末的 Darboux 前后吧。当时辛几何所涉及的，是分析力学即 Newton 的古典力学。

随着 20 世纪的到来，量子力学出现了。然而，从古典力学中诞生的辛几何，不仅没有结束其使命，反而与量子力学结下了更深、更神秘的关系。

为了作量子化，必须把古典理论用哈密顿形式来写。而且，上面强调的 $q$ 与 $p$ 的对等性，进入量子论后依然健在。也就是说，波函数既可以看作位置 $q$ 的函数，也可以看作动量 $p$ 的函数。两者之间的关系是 Fourier 变换。（顺带一提，波函数不能同时以 $p,q$ 两者为自变量，这正是不确定性原理的根据。）

然而，保持 $\omega$ 的一切变换即正则变换，是否都成为量子论的变换……

% [原文 p.13]
……并非如此。因此，并不存在从一般辛流形及其上的函数（哈密顿泛函）出发构造量子论的现成处方。

这一事实唤起了对前述立场的怀疑：“考虑辛流形，就是在非欧氏空间上展开分析力学”。这个立场原封不动地搬到量子论中是不成立的。

那么应当如何考虑呢？对此的解答至今仍未知晓。

这是笔者的想象：整体辛几何之所以没有顺利发展，理由之一或许就在这里。由于它过于紧密地与理论物理相关联而产生，许多辛几何专家或许会怀疑：若不摆脱对物理学出身的疑虑，纯粹作为数学来展开“通常的整体辛几何”，是否可行。\footnote{例如读 [296]，读者大概会惊讶于大森氏对流形所抱的、以近日几何学的状况来看堪称异样的否定性执著。笔者在那里看到了最近才起步的“辛拓扑”所无法穷尽的、辛几何的深根。然而，对于大森氏关于整体几何学的否定性见解（参见 [295] 的序文等），笔者并不赞同。}

在“Gromov 以前”的辛几何中曾盛加研究的若干要点，从“通常的整体几何学”的观点来看不易理解；只有立足于与物理学特别是量子论等的关联，其意义才显现出来——笔者有这种感觉。

作为这类问题的例子，下面说明可上溯到 Dirac 的几何量子化（geometric quantization）问题。

$X$ 上的辛结构 $\omega$ 如下导出 $C^\infty(M)$ 上的积结构 $\{\cdot,\cdot\}$。对函数 $f$，令由它生成的哈密顿向量场（Hamilton vector field）$X_f$ 为满足
\begin{equation}
  \omega(X_f,V)=df(V) \tag{1.2}
\end{equation}
对任意向量场 $V$ 都成立的唯一向量场。用它定义
\begin{equation}
  \{f,g\}=(df)(X_g)=X_g(f)=-X_f(g). \tag{1.3}
\end{equation}
% [Tutor Note: 原书此处写 $C^\infty(M)$，但本段讨论的空间记为 $X$，应为 $C^\infty(X)$（原书 p.13）。译文照原样保留。]

% [原文 p.14]
称 (1.3) 为泊松括号（Poisson bracket）。当 $\omega=\sum dq_i\wedge dp_i$ 时，
\begin{align}
  X_f &= \sum_i \frac{\partial f}{\partial p_i}\frac{\partial}{\partial q_i}
       -\sum_i \frac{\partial f}{\partial q_i}\frac{\partial}{\partial p_i}, \tag{1.4}\\
  \{f,g\} &= \sum_i \frac{\partial f}{\partial q_i}\frac{\partial g}{\partial p_i}
       -\sum_i \frac{\partial f}{\partial p_i}\frac{\partial g}{\partial q_i}. \tag{1.5}
\end{align}
泊松括号在 $\R$ 上双线性，且容易看出满足 Jacobi 律
\begin{equation}
  \{\{f,g\},h\}+\{\{g,h\},f\}+\{\{h,f\},g\}=0, \tag{1.6}
\end{equation}
也满足 $\{f,g\}=-\{g,f\}$。也就是说，泊松括号在 $C^\infty(M)$ 上定了李代数结构。（这个李代数通过对应 $f\mapsto X_f$，与保持 $X$ 的辛形式的微分同胚所成李群的李代数同构。参见 §2.2。）

设 $M=\R^{2n}$，其坐标为 $p_i,q_i$，$\omega=\sum dq_i\wedge dp_i$。把 $C^\infty(\R^{2n})$ 中关于诸 $p_i$ 为多项式的子集暂且记作 $A$。对 $A$ 的元素
\[
  F=\sum f_{\alpha_1,\dots,\alpha_n}(q_1,\dots,q_n)\,p_1^{\alpha_1}\cdots p_n^{\alpha_n},
\]
对应 $\R^{2n}$ 的微分算子
\begin{equation}
  \Omega(F)=\sum \sqrt{-1}^{\,\alpha_1+\cdots+\alpha_n}
  f_{\alpha_1,\dots,\alpha_n}(q_1,\dots,q_n)\,
  \frac{\partial^{\alpha_1+\cdots+\alpha_n}}{\partial q_1^{\alpha_1}\cdots\partial q_n^{\alpha_n}}. \tag{1.7}
\end{equation}
此时可以验证
\begin{equation}
  \Omega(\{F,G\})\equiv \sqrt{-1}\,[\Omega(F),\Omega(G)]
  \qquad \mathrm{mod}\ \deg F+\deg G-2\ \text{阶以下的项}. \tag{1.8}
\end{equation}
这里 $[\Omega(F),\Omega(G)]=\Omega(F)\Omega(G)-\Omega(G)\Omega(F)$ 是交换子。把这一点推广到未必是余切丛的辛流形上——这就是几何量子化的问题。也就是说：

\begin{problem}{1.5}
  对辛流形 $(M,\omega)$，构造 $C^\infty(M)$ 的尽可能大的子李代数 $A$、Hilbert 空间 $\mathcal{H}$，以及把 $A$ 的元素 $F$ 对应到线性算子 $\Omega(F)\colon\mathcal{H}\to\mathcal{H}$ 的对应 $\Omega$，使之满足 (1.8)。
\end{problem}

% [原文 p.15]
回想式 (1.7) 正是表达“把相空间上的函数看作算子”这一量子化手续的式子，问题 1.5 在量子力学中的意义就清楚了。此外，这个问题没有好的解答，正体现了前面所述的事实：把分析力学从余切丛推广到一般辛流形后，其量子化究竟是什么就变得不明确了。然而，离开量子力学、纯作为数学问题来看时，问题 1.5 的意义未必清楚。问题 1.5 将在 §4.3 中解说。

下面说明形变量子化（deformation quantization）的问题。

近来常听到非交换几何（noncommutative geometry）这个词。简单地说，非交换几何就是其上函数所成环为非交换的空间的几何学。

再稍作说明。对空间 $M$，考虑其上全体函数所成的环。说到“函数”，要考虑何种程度的函数，因情况而异。也就是说，若 $M$ 只是拓扑空间，则全体连续函数是合适的；若 $M$ 是微分流形，则全体光滑函数合适；若 $M$ 是复流形或代数簇，则全体全纯函数或多项式合适。无论如何，空间上的全体函数构成交换环。已知 $M$ 的几何学的许多部分，都可以翻译成这个环的性质。例如，$M$ 的点对应极大理想，$M$ 上的向量丛对应这个环上的模。（后者在考虑连续函数环 $C^0(M)$ 时更为明确，称为 Serre–Swan 定理。参见 [352] 等。）Gelfand 给出了使这个对应、即空间 $M\mapsto$ 环 $C^0(M)$ 更明确的如下定理。

\begin{theorem}{1.6}
  $M\mapsto C^0(M)$ 给出紧致 Hausdorff 空间全体与交换 $C^*$ 代数全体之间的 1 对 1 对应。
\end{theorem}

另一方面，对代数方程组
\[
  P_1(X_1,\dots,X_n)=\cdots=P_k(X_1,\dots,X_n)=0
\]
的零点集，代之以由 $P_1,\dots,P_k$ 生成的理想 $\mathfrak{m}$ 对多项式环 $\C[X_1,\dots,X_n]$ 的商环 $\C[X_1,\dots,X_n]/\mathfrak{m}$ 来考虑——这就是 Grothendieck……

% [原文 p.16]
……的概型（scheme）思想的出发点。

若把上述想法标语式地概括，就是：
\begin{center}
“空间”即“交换环”。
\end{center}
让我们试着把空间 $M$ 形变为与通常意义不同的“空间”。这里取 $M$ 为微分流形，考虑全体光滑函数所成的环 $C^\infty(M)$。按上面的标语，不必形变空间，只要形变环 $C^\infty(M)$ 即可。也就是说，尝试把交换环 $C^\infty(M)$ 形变为非交换环。把这个形变后的积记作 $f\star_\epsilon g$，即令
\begin{equation}
  f\star_\epsilon g=fg+\sum_k \epsilon^k B_k(f,g). \tag{1.9}
\end{equation}
这里 $fg$ 是普通的积。(1.9) 的右边暂且看作 $\epsilon$ 的形式幂级数。于是提出如下要求。

\begin{assumption}{1.7}
  $B_k(f,g)$ 关于 $f,g$ 双线性，并且关于 $f,g$ 是（有限阶的）微分算子。
\end{assumption}

也就是说，假定可表为
\begin{equation}
  B_k(f,g)=\sum_\ell\sum_{i_1,\dots,i_\ell,j_1,\dots,j_\ell}
  b_{i_1,\dots,i_\ell,j_1,\dots,j_\ell}\,
  \frac{\partial^\ell f}{\partial x_{i_1}\cdots\partial x_{i_\ell}}\,
  \frac{\partial^\ell g}{\partial x_{j_1}\cdots\partial x_{j_\ell}}. \tag{1.10}
\end{equation}
（这里 $b_{i_1,\dots,i_\ell,j_1,\dots,j_\ell}$ 是光滑函数。更准确地说，$b$ 是使 (1.10) 右边在坐标变换下不变、具有相应坐标变换性质的张量。）

再假定积 (1.9) 是结合的。也就是说，设
\begin{equation}
  (f\star_\epsilon g)\star_\epsilon h=f\star_\epsilon(g\star_\epsilon h) \tag{1.11}
\end{equation}
作为形式幂级数成立。这里定义 $\{f,g\}=B_1(f,g)$。由 (1.11) 可以验证 $\{\cdot,\cdot\}$ 满足 Jacobi 律 (1.6)。

\begin{definition}{1.8}
  $\{\cdot,\cdot\}\colon C^\infty\times C^\infty\to C^\infty$ 称为\textbf{泊松结构}（Poisson structure），是指满足以下三个条件：
  \begin{enumerate}[label=(\arabic*),leftmargin=2.4em,itemsep=1pt,topsep=2pt]
    \item $\{\cdot,\cdot\}$ 是 $\R$ 双线性的，且 $\{f,g\}=-\{g,f\}$；
    \item Jacobi 律 (1.6) 成立；
    \item $\{f,gh\}=g\{f,h\}+h\{f,g\}$。
  \end{enumerate}
\end{definition}

% [原文 p.17]
由 (3) 可知 $\{\cdot,\cdot\}$ 关于 $f,g$ 都是 1 阶双线性微分算子。辛结构由 (1.3) 确定一个泊松结构。一般的泊松结构可视为“退化的辛结构”。

\begin{remark}{1.9}
  在辛结构的情形，$\{\cdot,\cdot\}$ 的系数即 (1.10) 中的 $b_{ij}$，可由辛形式 $\sum\omega_{ij}dx_i\wedge dx_j$ 的系数 $\omega_{ij}$、利用 $\omega$ 自身作指标升降得到。（按张量分析中指标写法的约定，应当写作 $b^{ij}$ 而非 $b_{ij}$。）因此，退化的闭微分 2-形式与泊松结构是不同的。
\end{remark}

那么，形变量子化的问题指的是下述问题。

\begin{problem}{1.10}
  对给定的泊松结构 $\{f,g\}$，构造满足 $\{f,g\}=B_1(f,g)$ 的 $B_k(f,g)$，使之满足假设 1.7 与式 (1.11)。
\end{problem}

这个问题，在泊松结构 $\{\cdot,\cdot\}$ 由辛结构确定的情形，已由 De Wilde–Lecomte [70]、Omori–Maeda–Yoshioka [294]、Fedosov [99]、[100] 解决；在一般泊松结构的情形，则由 Kontsevich 解决 [221]（见 §6.5）。

% [Tutor Note: Kontsevich 的形量化定理把“泊松流形的形变量子化存在性”归结为 Hochschild 上同调与多向量场代数之间的 $L_\infty$ 拟同构（formality 定理）；这是 1990 年代数学物理最深刻的结果之一，本书 §6.4–§6.5 会回到它。]


\begin{remark}{1.11}
  当 $M=\C^n$、$\omega$ 为通常的辛形式时，$\star_\epsilon$ 曾被具体计算出来（见 §6.5）。又，令 $\epsilon=1$，则与 (1.7) 通过
  \[
    \Omega(f\star_1 g-g\star_1 f)=\sqrt{-1}\,[\Omega(f),\Omega(g)]
  \]
  相联系。
\end{remark}

刚才说明的两个问题，可认为是“Gromov 以前”辛几何的典型问题。（对于“Gromov 以前”就盛加研究的、与奇点论和微局部分析相关的部分，此处不再触及。）

当 Gromov 引入伪全纯曲线、“作为通常整体几何学”的辛几何开始发展之时，几何量子化与形变量子化的问题……

% [原文 p.18]
……与前者关系不大。而作为“通常的整体几何学”的辛几何，是在与数理物理无缘的情况下迈出步伐的。

然而，转机不久就到来了。1990 年代初，随着镜像对称（mirror symmetry）的发现，人们发现伪全纯曲线理论等价于在粒子物理中大行其道的超弦理论（super string theory）的一种，准确地说，是拓扑 $\sigma$ 模型（topological sigma model）。由伪全纯曲线确定的辛几何不变量，冠以 Gromov 与在物理侧引入拓扑 $\sigma$ 模型的 Witten 之名，被称为 Gromov–Witten 不变量；其研究成为镜像对称数学研究的中心。

而在其中，几何量子化与形变量子化都不知不觉地重新登场（见第 6 章）。然而，几何量子化和形变量子化为何会出现、原本关于量子化的问题意识与它们在镜像对称研究中的登场又有什么关联，坦率地说，笔者至今仍不清楚其“哲学上的道理”。当这一点明确之时，辛几何所怀有的对 21 世纪几何学的梦想，会连同整体问题一起，规模愈发宏伟地膨胀起来——我想。我相信那一天并不遥远。

辛几何另一重要侧面，即哈密顿动力系统的研究与伪全纯曲线的关系，是由 Floer 发现的。Floer 同调理论，是把伪全纯曲线应用于哈密顿动力系统周期轨道的研究（参见第 5 章）。然而，更具动力系统色彩的研究，例如 Kolmogorov–Arnold–Moser 理论（KAM 理论）与混沌等和伪全纯曲线的关系，目前似乎并不十分清楚\footnote{§2.7 所述的应用，是否与此相关呢。}。

另一方面，与辛几何相关的另一重要领域——奇点论——与伪全纯曲线的关系，从镜像对称发现之初就十分显著……

% [原文 p.19]
……十分显著。例如，斋藤恭司的原始形式（primitive form）与平坦坐标（flat coordinate）理论（虽属奇点论），在镜像对称研究中发挥了重要作用（关于这一点参见 [321]、[322]、[243] 等）。此外，振荡积分以及出现在奇点附近的接触流形等，也在 Gromov–Witten 不变量研究中发挥了重要作用（参见 [94]、[147] 等）。

% [Tutor Note: 斋藤恭司（Kyoji Saito）的“原始形式”理论，是他在 1970 年代为研究奇点的局部环与平坦坐标而建立的；本书 §6 会看到它如何成为 B 模型与镜像对称的核心工具。文献号 [321] 等对应该书参考文献表。]

\chapterreview{本章是全书的问题意识与路线图。第 1 节用埃尔朗根纲领与 Cartan 的 G 结构把“几何”统一为群作用下的不变量理论，并指出辛结构正是 $G=Sp(n)\subset GL(2n;\R)$ 的 G 结构、其可积性条件为 $d\omega=0$——这为第 2 章“辛流形＝$(X,\omega)$”的定义作了观念上的铺垫。第 2 节以 1950 年代微分拓扑学与复几何的独立为起点，勾出“通常的整体几何学”这一参照系：几何学的中心从局部不变量转向整体存在性与分类问题；Gromov 1985 年的伪全纯曲线方法正是把这一范式带入辛几何的转机。第 3 节回到 Hamilton 力学，说明 $q$ 与 $p$ 的对等性如何把辛形式推到舞台中央，并借几何量子化（问题 1.5）与形变量子化（问题 1.10）两个“Gromov 以前”的典型问题，交代辛几何与量子论之间未了的瓜葛。作者反复申明本章是“正史”而非全史：辛几何长期被“量子化”这一梦想压着走，缺少像样的大定理；他期待 21 世纪的几何学会把正史之外的潮流推到前台。读本章时宜把它当作全书的“为什么”：后续第 2 章给工具（伪全纯曲线与模空间），第 5、6 章才回收本章许下的两个承诺（Floer 同调与镜像对称）。注意 §1.1 的 G 结构语言与 §2.3 的近复结构语言是同一件事的两面，初读时容易把它们当成两套独立的形式体系。}
